\section{Capítulo~\ref{sec:neurontypes}. Tipos de neuronas}

Las cifras del capítulo se apoyan en recuentos directos de células en el cerebro humano allí donde existen; la regla “una neurona, un neurotransmisor”, en atlas celulares y en experimentos que también encontraron excepciones; el papel de \nt{DOPA}, en registros de espigas; y los papeles de los demás canales de difusión, en hipótesis.

{\footnotesize
\setlength{\tabcolsep}{3pt}\renewcommand{\arraystretch}{1.15}
\begin{longtable}{>{\raggedright}p{0.5cm}>{\raggedright}p{4.0cm}>{\raggedright}p{5.6cm}>{\raggedright}p{2.3cm}>{\raggedright\arraybackslash}p{1.7cm}}
\toprule
N.º & Afirmación & Experimento y qué se midió & Fuente & Estado \\
\midrule\endfirsthead
\toprule
N.º & Afirmación & Experimento y qué se midió & Fuente & Estado \\
\midrule\endhead
1 & El cerebro humano tiene unas $8.6\cdot10^{10}$ neuronas, y casi todas las neuronas \nt{GLUT} son pequeñas neuronas del cerebelo (tabla~\ref{tab:types}) & Fraccionador isotrópico: el tejido se disuelve en una suspensión homogénea de núcleos y se cuentan los núcleos con el marcador neuronal NeuN. Los varones adultos tienen $86.1\pm8.1$ mil millones de neuronas; en la corteza cerebral está solo el $19\%$, y la mayor parte, en el cerebelo & \cite{Azevedo2009} & medido \\
2 & Casi toda neurona tiene un neurotransmisor principal (proposición~\ref{prop:dale}), pero hay excepciones & Atlas transcriptómico del cerebro del ratón: unos $4\cdot10^6$ células, $5322$ grupos; a cada grupo se le asigna un neurotransmisor por los genes de síntesis y empaquetado. Las excepciones se midieron directamente: las neuronas \nt{DOPA} en rebanadas liberan también GABA a través del mismo transportador vesicular e inhiben neuronas del estriado; en parte de los axones \nt{DOPA}, el glutamato y la dopamina salen de terminales distintas del mismo cable (microscopía electrónica y optogenética) & \cite{Strata1999,Yao2023,Tritsch2012,Zhang2015} & medido \\
3 & Toda la columna de la matriz de pesos tiene un solo signo~\eqref{eq:dalesign} & Deducción del capítulo a partir de la proposición~\ref{prop:dale} y de que casi todos los receptores de glutamato son excitadores y los de GABA, inhibidores. No hay una comprobación directa de “todos los destinatarios de una neurona reciben un peso del mismo signo” como regla general; contraejemplos: la coliberación de GABA por neuronas \nt{DOPA} (fila~2) y los destinatarios con receptores de signo distinto (fila~11) & deducción del capítulo & teoría \\
4 & Entre tres cuartas y cuatro quintas partes de las neuronas de la corteza son \nt{GLUT}, entre una quinta y una cuarta parte, \nt{GABA} & Inmunotinción de GABA en columnas de $50$\,\textmu m de ancho a través de todo el espesor cortical en $10$ áreas de la corteza de mono: en $8$ áreas las neuronas GABA son cerca del $25\%$ de todas las neuronas; en la corteza visual primaria, menos del $20\%$ & \cite{Hendry1987} & medido \\
5 & Una neurona \nt{GLUT} tiene unas $10^4$ salidas & Estereología post mortem: en la neocorteza de varones jóvenes hay $1.64\cdot10^{14}$ sinapsis (microscopía electrónica, disector) y unas $2.3\cdot10^{10}$ neuronas. El cociente da $\sim7\cdot10^3$ sinapsis por neurona; en promedio, el número de entradas es igual al de salidas & \cite{Tang2001synapses,Pakkenberg1997} & medido \\
6 & La salida de una neurona \nt{DOPA} se ramifica en $10^5$--$10^6$ terminales; es una estimación por cobertura de la diana, no un recuento & Marcaje viral de la membrana de neuronas \nt{DOPA} individuales de rata y reconstrucción completa del axón: una neurona cubre entre el $0.45$ y el $5.7\%$ (en promedio el $2.7\%$) del volumen del estriado. Se midió la cobertura, no el número de terminales: este se dedujo de la cobertura en orden de magnitud, sin recuento directo. Para \nt{SERO} y \nt{NORA}, los números de terminales son estimaciones burdas & \cite{Matsuda2009} & medido \\
7 & Las neuronas de difusión son pocas: \nt{DOPA} $\sim5\cdot10^5$ (el principal de dos grupos, cota inferior), \nt{SERO} $\sim3\cdot10^5$ (suma de dos recuentos parciales), \nt{HIST} $\sim6\cdot10^4$ (tabla~\ref{tab:types}, fig.~\ref{fig:typepie}) & Recuentos estereológicos en humanos: $550\,000$ neuronas pigmentadas en la sustancia negra (sin el área tegmental ventral); $235\,000$ neuronas en el núcleo dorsal del rafe (todas las del núcleo) y otras $125\,000$ neuronas serotoninérgicas en el tegmento pontino; unas $64\,000$ neuronas histaminérgicas del hipotálamo. Para \nt{NORA}, \nt{GLYC}, \nt{ACET} y \nt{ADRE} no hay recuento directo: en la tabla son estimaciones burdas del orden de magnitud & \cite{Pakkenberg1991,Baker1990,Baker1991,Airaksinen1991} & medido \\
8 & El glutamato en la hendidura sube hasta cerca de un milimolar y decae en $\sim1\ms$ & Por el desplazamiento de un bloqueador competitivo de disociación rápida de los receptores NMDA durante la transmisión sináptica (cultivo de neuronas de hipocampo): pico de $1.1$\,mM, constante de decaimiento de $1.2\ms$ & \cite{Clements1992} & medido \\
9 & En la corteza activa, las corrientes excitadora e inhibidora están casi equilibradas y la neurona se mantiene cerca del umbral & Teoría: una red con conexiones fuertes llega por sí sola a un régimen equilibrado. Experimento: en la corteza prefrontal del hurón in vivo, durante los estados “up”, el potencial de inversión de la corriente sináptica se mantiene cerca de $-37$\,mV durante cientos de milisegundos, aunque la conductancia cambia en $\sim21$\,nS: la excitación y la inhibición suben y bajan juntas & \cite{vanVreeswijk1996,Haider2006} & medido \\
10 & Las neuronas \nt{DOPA} comunican el error de predicción de la recompensa~\eqref{eq:rpe} (fig.~\ref{fig:rpe}) & Espigas de neuronas \nt{DOPA} de mono: ráfaga ante el jugo inesperado; tras el aprendizaje, ráfaga ante la señal y ninguna respuesta al jugo predicho; caída de la frecuencia cuando no se da el jugo prometido. En el experimento cuantitativo, la frecuencia es proporcional a la diferencia entre la recompensa actual y una media ponderada de las pasadas, pero solo para errores positivos. La ecuación~\eqref{eq:rpe} misma es el algoritmo de diferencias temporales & \cite{Schultz1997,Bayer2005,SuttonBarto2018} & medido \\
11 & El signo y la velocidad los fija el receptor: los ionotrópicos, fracciones de milisegundo; los metabotrópicos, mucho más lentos (proposición~\ref{prop:receptor}) & Pulsos rápidos de glutamato sobre parches de membrana de neuronas de hipocampo: la corriente por los receptores ionotrópicos sube (del $20$ al $80\%$) en $0.2$--$0.6\ms$ y decae en $2$--$3\ms$. El potencial inhibidor lento de las neuronas piramidales se suprime con un bloqueador selectivo del receptor metabotrópico de GABA. Dos familias de receptores de dopamina con acción opuesta; en el estriado, las neuronas con receptor D1 necesitan dopamina para reforzar sinapsis, y las neuronas con D2, para debilitarlas & \cite{Colquhoun1992,DutarNicoll1988,Beaulieu2011,Shen2008} & medido \\
12 & Un canal de difusión no es un solo cable, sino un mazo de pocas líneas (sección~\ref{sec:buses}) & Marcaje genético-viral de las neuronas serotoninérgicas del núcleo dorsal del rafe del ratón: las neuronas que envían su salida a la amígdala y a la corteza frontal ocupan lugares distintos del núcleo, se ramifican en regiones complementarias y responden de forma opuesta a un estímulo desagradable. Revisión: subgrupos de neuronas del locus coeruleus (\nt{NORA}) codifican procesos distintos & \cite{Ren2018,Poe2020} & medido \\
13 & \nt{NORA} fija la ganancia $g$ & Hipótesis de la ganancia adaptativa; modelo de red en el que las catecolaminas aumentan la pendiente de la característica de transferencia. No hay una medición directa de “$g$ crece con la noradrenalina” en toda la red; se midió que el diámetro de la pupila sigue las espigas de las neuronas del locus coeruleus del mono & \cite{AstonJones2005,ServanSchreiber1990,Joshi2016} & hipótesis \\
14 & \nt{ACET} conmuta “escritura/lectura” y fija la tasa de aprendizaje & Hipótesis. Apoyo: en rebanadas de corteza olfativa de rata, los agonistas de la acetilcolina ($100$\,\textmu M) debilitan las sinapsis internas, las que “recuerdan”, en más del $60\%$, y las de entrada, en menos del $15\%$ & \cite{Hasselmo2006,HasselmoBower1992} & hipótesis \\
15 & \nt{SERO} fija el factor de descuento $\gamma$; las neuronas serotoninérgicas trabajan de forma regular, unas pocas espigas por segundo, y casi callan en una de las fases del sueño & Hipótesis “serotonina $=\gamma$”. El argumento más fuerte: la activación optogenética de las neuronas serotoninérgicas del núcleo dorsal del rafe del ratón reduce el número de abandonos prematuros de la espera de una recompensa diferida. La frecuencia de espigas y el silencio durante el sueño de movimientos oculares rápidos están medidos (revisión de registros) & \cite{Doya2002,Miyazaki2014,JacobsAzmitia1992} & hipótesis \\
\bottomrule
\end{longtable}}
